\subsection{The Ember Days of Advent}\label{sec:ember}

The Ember Days of December are the Wednesday, Friday and Saturday after the Third Sunday of
Advent. In 2026 they fall on 16, 18 and 19 December. They keep the order St Leo the Great gave his people
for the fast of the tenth month: fast on Wednesday and Friday, and keep vigil on Saturday with the
Apostle Peter (Section~\ref{sec:history}). In the Breviary of 1888, which prints Leo's sermon at
Matins of the Third Sunday, the notice reads:

\begin{quote}
\latin{Quarta igitur et sexta Feria jejunemus: Sabbato autem apud beatum Petrum Apostolum pariter
vigilemus:}\par
{\smallskip\noindent\small ---Leo the Great, \textit{Sermo} 13, in \textit{Breviarium Romanum},
pars hiemalis (Ratisbon: Pustet, 1888), pp.~153--154. Gloss: on Wednesday and Friday, then, let us
fast, and on Saturday let us keep vigil together with blessed Peter the Apostle.\par}
\end{quote}

\subsubsection{The three days}

On the Wednesday the station is at St Mary Major. The Mass reads two lessons of Isaiah (Is 2:2--5;
7:10--15) and the Gospel of the Annunciation, \latin{Missus est} (Luke 1:26--38). The second
lesson is the sign given to Ahaz:

\begin{quote}
Therefore the Lord himself shall give you a sign. Behold a virgin shall conceive, and bear a son
and his name shall be called Emmanuel.\par
{\smallskip\noindent\small ---Isaias 7:14, Douay-Rheims (Challoner).\par}
\end{quote}

\noindent On the Friday the station is at the Twelve Apostles. The Mass reads Isaiah's rod out of
the root of Jesse (Is 11:1--5; Section~\ref{sec:second-sunday}) and the Gospel of the Visitation
(Luke 1:39--47), and its Collect is the third of the season's prayers that begin
\latin{Excita}. On the Saturday the station is at St Peter's. In the longer form of its Mass,
which the Missal appoints for the conventual Mass and for a Mass at which orders are conferred,
the last of four lessons of Isaiah is the oracle to Cyrus that ends with the verse the Fourth
Sunday takes up (Is 45:1--8; Section~\ref{sec:fourth-sunday}); the Gospel is John's preaching in
the desert (Luke 3:1--6).

At Matins the Ember days read no Isaiah. Each has three lessons of homily on its Gospel: St Ambrose
on the Wednesday and the Friday, St Gregory on the Saturday (n.~222\,a). Ambrose begins the
Wednesday homily by confessing that the counsels of God are unsearchable, and then asks why she
who was chosen to be the Mother of the Lord was espoused to a man before the power of the Most
High overshadowed her. On the Friday he follows Mary into the hill country:

\begin{quote}
She went joyfully as one who hath received a mercy in answer to his vow goeth to pay the same. She
went with devotion, as a godly person goeth to execute a religious duty. She went into the hill
country in joyful haste. And is it not something that she went up into the hills? God was already
in her womb, and her feeling bore her continually upward. The grace of the Holy Spirit knoweth no
slow working.\par
{\smallskip\noindent\small ---Ambrose, \textit{Expositio evangelii secundum Lucam} II.19, tr.
Bute, \textit{The Roman Breviary} (1908), Winter, pp.~252--253.\par}
\end{quote}

\noindent The last sentence in Ambrose's Latin reads \latin{nescit tarda molimina sancti spiritus
gratia} (CSEL 32/4, p.~52). The Saturday homily is Gregory's twentieth, which he preached on an
Ember Saturday (Section~\ref{sec:history}); it opens with the rulers' names by which the Gospel
sets the Baptist's preaching in history.

\subsubsection{The world waits on her answer}

St Bernard, preaching on the Wednesday's Gospel, stands at the moment before Mary's answer and
speaks to her for the whole race:

\begin{quote}
Behold the Angel now awaits thy answer: ``it is time that he should return to the Lord Who sent
him.'' We also, O Lady, await from thy lips the sentence of mercy and compassion, we who are so
miserably groaning under the sentence of condemnation. For lo! the price of our salvation is now
offered to thee: if thou wilt only consent, we shall at once be set at liberty. We have been
created by the eternal Word of God, and behold we die: by thy momentary word we must be renewed
and restored to life.\par
{\smallskip\noindent\small ---Bernard, \textit{Super Missus est} IV.8, tr. a priest of Mount
Melleray (Dublin, [1921]).\par}
\end{quote}

A little further on Bernard presses her: ``Say the word and receive the Word.'' And again:
``Arise by faith, make haste by devotion, open by consent'' (\textit{Super Missus est} IV.8, tr.
Mount Melleray). The Gospel gives her answer:

\begin{quote}
And Mary said: Behold the handmaid of the Lord: be it done to me according to thy word. And the
angel departed from her.\par
{\smallskip\noindent\small ---Luke 1:38, Douay-Rheims (Challoner).\par}
\end{quote}

\noindent The Annunciation on Wednesday gives her consent, and the Visitation on Friday her haste.

\subsubsection{Fasting, prayer and alms}

Leo preached the December fast as one exercise of three. In one sermon on the fast he joins it to
prayer and almsgiving:

\begin{quote}
For by prayer we seek to propitiate God, by fasting we extinguish the lusts of the flesh, by alms
we redeem our sins: and at the same time God's image is throughout renewed in us, if we are always
ready to praise Him, unfailingly intent on our purification and unceasingly active in cherishing
our neighbour.\par
{\smallskip\noindent\small ---Leo the Great, \textit{Sermo} XII.4, tr. C.~L. Feltoe, NPNF, 2nd
ser., vol.~12 (New York, 1895).\par}
\end{quote}

In the sermon whose closing notice is quoted above, which the Breviary read at Matins of the Third
Sunday before 1960, he puts the same teaching in a few short sentences:

\begin{quote}
Fasting hath ever been the bread of strength. From abstinence proceed pure thoughts, reasonable
desires, and healthy counsels. By voluntary mortifications the flesh dieth to lust, and the soul
is renewed in might. But since fasting is not the only mean whereby we get health for our souls,
let us add to our fasting works of mercy. Let us spend in good deeds what we take from indulgence.
Let our fast become the banquet of the poor.\par
{\smallskip\noindent\small ---Leo the Great, \textit{Sermo} 13, tr. Bute, \textit{The Roman
Breviary} (1908), Winter, Third Sunday of Advent, lesson v.\par}
\end{quote}

Under the Code of Canon Law of 1917, in force for the Latin Church from 1918 to 1983, the Ember
days of December were days of fast and abstinence (can.~1252 \S\,2). Neither the Code of 1983 nor,
in the United States, the bishops' norms of 1966, both in force in 2026, binds anyone to fast on
them.
Section~\ref{sec:fasting} sets out the law and the custom in full.

\medskip
\noindent\textit{Counsel for the week.} If you can, hear Mass on one of the three days, and read
its Gospel beforehand: the Annunciation on Wednesday, the Visitation on Friday. Keep the three
days with some fast of your own choosing, within what your health and duties allow and with your
confessor's advice if you have one, and give to the poor what the fast saves, as Leo asks. On the
Wednesday read Bernard's appeal to the Virgin slowly and make his last words your own prayer. The
Missal provides for orders to be conferred at the Saturday Mass; pray that day for those who will be
ordained.
